% from ega2-1-fr.tex
\oldpage[II]{81}

$Y$-immersion $r:X\to P=\mathbf{P}(\mathcal{E})$ telle que
$\mathcal{L}=r^*(\mathcal{O}_P(1))$. Posons
$Q=\mathbf{P}(\mathcal{E}\otimes\mathcal{E}')$ et considérons le morphisme de
Segre $\varsigma:P\times_YP'\to Q$ (\hyperref[II.4.3.1-fr]{4.3.1}). Comme
$r$ est une immersion, il en est de même de
$(r,r')_Y:X\to P\times_YP'$ (\hyperref[I.5.3.14-fr]{I, 5.3.14})~; mais
$\varsigma$ étant une immersion (\hyperref[II.4.3.3-fr]{4.3.3}), il en est de
même de
$r'':X\xrightarrow{(r,r')_Y}P\times_YP'\xrightarrow{\varsigma}Q$. D'autre
part (\hyperref[II.4.3.2.1-fr]{4.3.2.1}),
$\varsigma^*(\mathcal{O}_Q(1))$ est isomorphe à
$\mathcal{O}_P(1)\otimes_Y\mathcal{O}_{P'}(1)$, donc
(\hyperref[I.9.1.4-fr]{I, 9.1.4}), $r''^*(\mathcal{O}_Q(1))$ est isomorphe à
$\mathcal{L}\otimes\mathcal{L}'$, ce qui démontre la proposition.
\end{proof}

\begin{corollary}[4.4.9]
\label{II.4.4.9-fr}
Soit $q:X\to Y$ un morphisme.
\begin{enumerate}
  \item[\rm{(i)}] Soient $\mathcal{L}$ un $\mathcal{O}_X$-Module inversible,
  $\mathcal{K}$ un $\mathcal{O}_Y$-Module inversible. Pour que $\mathcal{L}$
  soit très ample relativement à $q$, il faut et il suffit que
  $\mathcal{L}\otimes q^*(\mathcal{K})$ le soit.
  \item[\rm{(ii)}] Si $\mathcal{L}$ et $\mathcal{L}'$ sont deux
  $\mathcal{O}_X$-Modules très amples relativement à $q$, il en est de même
  de $\mathcal{L}\otimes\mathcal{L}'$~; en particulier
  $\mathcal{L}^{\otimes n}$ est très ample relativement à $q$ pour tout
  $n>0$.
\end{enumerate}
\end{corollary}

\begin{proof}
(ii) est conséquence immédiate de (\hyperref[II.4.4.8-fr]{4.4.8}), ainsi que
la nécessité de la condition (i)~; d'autre part, si
$\mathcal{L}\otimes q^*(\mathcal{K})$ est très ample, il en est de même de
$(\mathcal{L}\otimes q^*(\mathcal{K}))\otimes q^*(\mathcal{K}^{-1})$
d'après ce qui précède, et ce dernier $\mathcal{O}_X$-Module est isomorphe à
$\mathcal{L}$ (\hyperref[0.5.4.3-fr]{0, 5.4.3} et
\hyperref[0.5.4.5-fr]{5.4.5}).
\end{proof}

\begin{proposition}[4.4.10]
\label{II.4.4.10-fr}
\begin{enumerate}
  \item[\rm{(i)}] Pour tout préschéma $Y$, tout $\mathcal{O}_Y$-Module
  inversible $\mathcal{L}$ est très ample relativement au morphisme identique
  $1_Y$.
  \item[\rm{(i bis)}] Soient $f:X\to Y$ un morphisme, $j:X'\to X$ une
  immersion. Si $\mathcal{L}$ est un $\mathcal{O}_X$-Module très ample
  relativement à $f$, $j^*(\mathcal{L})$ est très ample relativement à
  $f\circ j$.
  \item[\rm{(ii)}] Soient $Z$ un préschéma quasi-compact, $f:X\to Y$ un
  morphisme de type fini, $g:Y\to Z$ un morphisme quasi-compact,
  $\mathcal{L}$ un $\mathcal{O}_X$-Module très ample relativement à $f$,
  $\mathcal{K}$ un $\mathcal{O}_Y$-Module très ample relativement à $g$.
  Alors il existe un entier $n_0>0$ tel que
  $\mathcal{L}\otimes f^*(\mathcal{K}^{\otimes n})$ soit très ample
  relativement à $g\circ f$ pour tout $n\geq n_0$.
  \item[\rm{(iii)}] Soient $f:X\to Y$, $g:Y'\to Y$ deux morphismes, et posons
  $X'=X_{(Y')}$. Si $\mathcal{L}$ est un $\mathcal{O}_X$-Module très ample
  relativement à $f$, $\mathcal{L}'=\mathcal{L}\otimes_Y\mathcal{O}_{Y'}$
  est un $\mathcal{O}_{X'}$-Module très ample relativement à $f_{(Y')}$.
  \item[\rm{(iv)}] Soient $f_i:X_i\to Y_i$ ($i=1,2$) deux $S$-morphismes.
  Si $\mathcal{L}_i$ est un $\mathcal{O}_{X_i}$-Module très ample relativement
  à $f_i$ ($i=1,2$), $\mathcal{L}_1\otimes_S\mathcal{L}_2$ est très ample
  relativement à $f_1\times_Sf_2$.
  \item[\rm{(v)}] Soient $f:X\to Y$, $g:Y\to Z$ deux morphismes. Si un
  $\mathcal{O}_X$-Module $\mathcal{L}$ est très ample relativement à
  $g\circ f$, alors $\mathcal{L}$ est très ample relativement à $f$.
  \item[\rm{(vi)}] Soient $f:X\to Y$ un morphisme, $j$ l'injection canonique
  $X_{\mathrm{red}}\to X$. Si $\mathcal{L}$ est un $\mathcal{O}_X$-Module
  très ample relativement à $f$, alors $j^*(\mathcal{L})$ est très ample
  relativement à $f_{\mathrm{red}}$.
\end{enumerate}
\end{proposition}

\begin{proof}
La propriété (i bis) découle aussitôt de la définition
(\hyperref[II.4.4.2-fr]{4.4.2}) et il est immédiat que (vi) se déduit
formellement de (i bis) et de (v) par un raisonnement calqué sur celui de
(\hyperref[I.5.5.12-fr]{I, 5.5.12}), que nous laissons au lecteur. Pour
démontrer (v), considérons, comme dans
(\hyperref[I.5.5.12-fr]{I, 5.5.12}), la factorisation
$X\xrightarrow{\Gamma_f}X\times_ZY\xrightarrow{p_2}Y$, en notant que
$p_2=(g\circ f)\times1_Y$. Il résulte de l'hypothèse et de (i) et (iv) que
$\mathcal{L}\otimes_{\mathcal{O}_Z}\mathcal{O}_Y$ est très ample pour $p_2$~;
d'autre part, on a
$\mathcal{L}=\Gamma_f^*(\mathcal{L}\otimes_{\mathcal{O}_Z}\mathcal{O}_Y)$
(\hyperref[I.9.1.4-fr]{I, 9.1.4}), et $\Gamma_f$ est une immersion
(\hyperref[I.5.3.11-fr]{I, 5.3.11})~; on peut donc appliquer (i bis).

